De sluwe charlatan
toonde een jongeman,

die onschuld uitstraalde,
maar doelloos ronddwaalde,

trots zijn VOC-schip.
Die keek eerst blij, toen sip.

Zo’n vijftig meter lang,
zo’n vijf meter diepgang.

Drie masten stonden fier,
de dekken telden vier.

De romp van eikenhout
was uit planken gebouwd

op gebogen spanten.
En van alle kanten

zag men versierselen
met zeese schepselen.

Het leek hem fantastisch
om op zo’n bombastisch

vaartuig mee te varen.
Hij wou niet verklaren

waarom hij bedroefd was,
als ook wel in zijn sas.

Hij kwam uit de Jordaan
maar monsterde niet aan.

Hoewel hij kansrijk leek
zo komend uit de streek.

De profeet dacht toen: "bah."
Het was mauvaise foi,

zei ooit een filosoof,
oftewel slecht geloof.

Hij verschool zich pardoes
achter een valse smoes.

Leve de prins-bisschop
was niet het vervolg op

de prins-bisschop is dood.
Het bisdom was in nood

en zocht een vervanger.
Dit keer geen aanhanger

van de Zonnekoning,
liefst ook geen ouderling.

Niet de hypochondrie
of de melancholie,

die de vrome einstein
tijdens zijn ambtstermijn

chronisch etaleerde
en zo escaleerde

met menig onderdaan,
die hem het liefst zag gaan.

Zelfs de God van het licht
was op den duur gezwicht

voor Loki’s sluwe list
door een pijl die niet mist.

Ook Balder bleek kwetsbaar
en tevens vervangbaar

net als de oude bisschop.
Het kostte hem zijn kop.

De profeet als kwaadste
bezocht hem als laatste.

Hij was terechtgesteld
en leek teleurgesteld

in zijn presentatie
bij de convocatie

die de prins-bisschop hield.
Zijn leer werd daar vernield.

De profeet dramde door.
Zo kwam hem ten gehoor

dat vast de Jordanees,
die zijn retourschip prees,

zou worden verweten
te hebben gezeten,

voor langer dan vijf jaar
- met berouw weliswaar -

en daardoor geen kans had.
Sollicitanten zat

met een brandschone lei.
“Die kiest men boven mij!

Zeker in vredestijd.”
Het formele beleid

sloot op dit verhaal aan.
Daar kon men omheen gaan.

De profeet was verrast
dat zijn vredige gast

crimineel leek te zijn.
Misschien was zijn faam schijn.

Om toen in de Jordaan
rechtop te blijven staan

bedienden velen zich
van een buitenissig

en strafbaar lulverhaal.
Het draaide allemaal

om het erbij horen.
Met dovemansoren

vroeg de profeet stug door.
“Waar zat je dan ooit voor?”

“Ik stal een stukje brood!”
“Ik ben geen idioot

die gelooft dat die straf
hoort bij wat je aangaf.

Want bij zulke streken
hoort dagen of weken.”

“Nochtans zei de vierschaar
dat eens het roven maar

afgelopen moest wezen.”
Men bleek hem te vrezen

voor het harde geweld
waarmee men werd geveld.

Een gerimpelde hand
trilde aan elke kant

hangend onder tijdsdruk
vlak boven een schaakstuk.

Er was nog geen schaakklok,
doch het zwijgend gemok,

geuit via gapen,
wist hem toch te kapen.

Die sfeer van ongeduld
gaf zijn sloomheid de schuld.

Hij vermoedde een flop,
maar tilde het paard op,

omwille van zijn trots.
Hij besefte toen plots

dat hij had geblunderd.
Zijn paard werd geplunderd.

De geschokte bisschop
gaf daarna meteen op.

Toen een andere hand
vanuit de overkant

zijn koningin bewoog,
keek de bisschop omhoog.

Hij miste een schaakmat,
die hij kennelijk had

en aan hem werd getoond.
Zijn gehaast werd gehoond.

Schaken was zijn training
in zelfontwikkeling.

Om de geest te scherpen
en te onderwerpen

verlies en frustratie.
Hij moest imperfectie,

ondanks zijn intellect,
aanvaarden als aspect

van zijn aardse bestaan.
Ermee leren omgaan.

We zijn een mengeling
van rede én neiging,

zei Immanuel Kant.
Een tweeslachtig verband.

Zijn beste schaakmaatje
werd zijn kameraadje;

de landheer van Emmy.
Zijn speldominantie

kwam niet vanuit techniek
of hersengymnastiek

maar uit psychologie.
Met een blufsuggestie,
zoals een snelle zet,
wierp hij zijn psychisch net.

Niet dat het uitmaakte.
Want het koppel schaakte

puur ter relaxatie
voor de convocatie.

““Immers geen emotie,
die mag zijn, blijft”, zegt ie.

Dat is Oosters gezwets.
Wat ik u als raad schets

is Psalm 19:10,
waarin u dat kunt zien

en dit feit kan leren;
De vreze des Heeren

is niet alleen reinheid,
maar voor de eeuwigheid.”

De bisschop telkens weer
ontkrachtigde zijn leer

en de kerkraad volgde.
De profeet vervolgde

zijn leerverdediging.
“Nee zeg ik. Gods schepping

is niet problematisch,
maar probabilistisch!”

De bisschop rook zijn kans,
maar behield zijn balans.

“Oh, mijn handen jeuken.
Men leest toch in Spreuken

dat de Heer door wijsheid
de aarde heeft gekleid.”

Alle aanwezigen
konden dit bezigen

en knikten instemmend.
Het werkte ontremmend.

“De bijbel is verjaard
als men het diep aankaart.

Maak geen blinde knieval.
Kijk eens naar het heelal.

Hoe meer men er naar kijkt,
hoe ouder de zaak lijkt!”

De bisschop won weer fier;
“Lees Korintiërs vier.

Met name vers achttien.
Daar zult u zijn waan zien.

Want hij heeft het weer mis.
“Want wat onzichtbaar is

is eeuwig” leest men daar.
Zijn stelling is niet waar!”

De profeet zonder staf
beet gekwetst van zich af.

“Dat een klein brein klein denkt
en de waarheid daardoor krenkt,

is niet te voorkomen.
Het verstand intomen,

ondanks een begaafd brein,
bezorgt een mens hoofdpijn.”

Hij goot de bisschopswijn,
waar iets van het venijn

in het glas achterbleef,
daar hij zijn wraak overdreef,

opgejaagd in een hoek.
Hij nam uit een bidboek,

dat op een tafel lag,
een open briefomslag,

gericht aan de kerkraad
met de bisschops mandaat,

en gooide uit censuur
zijn doem in het haardvuur.

“Vergeeft u mij, eerwaarde
dat ik u niet spaarde?

In mijn goddelijk plan
past niet uw hersenpan,

wiens inhoud ik hoog acht.
Dat had ik niet verwacht.

Mijn zelfzekerheid
werd snel onwetendheid.

Heb uw vijanden lief,
maar blijf competitief.”

De profeet riep de wacht.
Hij werd geenszins verdacht,

want niemand was verward
wegens zijn zwakke hart.

Ook de aanklacht verdween.
De raad ging nietsdoend heen.
